\documentclass[11pt]{article}
\usepackage[margin=1in]{geometry}
\usepackage{amsmath, amssymb, amsthm}
\usepackage{hyperref}

\title{Math - Test Sections Questions: Change of Basis}
\author{Generated Problem Set}
\date{}

\begin{document}

\maketitle
\tableofcontents
\newpage

\section{Concept Checklist \& Question Mapping}
Based on the provided document ``Homework 4.6: Change of Basis'', the following concepts, problem types, tricks, and pitfalls have been identified. The 30 problems in this document are mapped to these concepts to ensure comprehensive coverage.

\begin{itemize}
    \item \textbf{C1: Direct Matrix Construction:} Extracting coefficients directly from linear combinations to form $P_{\mathcal{C} \leftarrow \mathcal{B}}$. (Questions: 1, 2, 30)
    \item \textbf{C2: Row Reduction Method in $\mathbb{R}^n$:} Using the augmented matrix $[\mathcal{C} \mid \mathcal{B}] \sim [I \mid P_{\mathcal{C} \leftarrow \mathcal{B}}]$. (Questions: 3, 4, 6, 8, 28)
    \item \textbf{C3: Coordinate Transformation:} Using the equation $[x]_{\mathcal{C}} = P_{\mathcal{C} \leftarrow \mathcal{B}}[x]_{\mathcal{B}}$. (Questions: 2, 4, 11, 30)
    \item \textbf{C4: Invertibility \& Reversing Direction:} Using $P_{\mathcal{B} \leftarrow \mathcal{C}} = (P_{\mathcal{C} \leftarrow \mathcal{B}})^{-1}$. (Questions: 5, 7, 29)
    \item \textbf{C5: Abstract Vector Spaces (Polynomials):} Isomorphic mapping of $\mathbb{P}_n$ to $\mathbb{R}^{n+1}$. (Questions: 9, 10, 11, 12)
    \item \textbf{C6: Abstract Vector Spaces (Functions/Trig):} Mapping continuous functions to discrete matrices. (Questions: 13, 14)
    \item \textbf{C7: Conceptual True/False:} Testing definitions, linear independence, and matrix properties. (Questions: 15, 16, 17, 18)
    \item \textbf{C8: Theoretical Proofs \& Derivations:} Relying on definitions of linear transformations. (Questions: 19, 20)
    \item \textbf{C9: Pitfall - Reversing the Direction:} Setting up the augmented matrix backwards as $[\mathcal{B} \mid \mathcal{C}]$. (Question: 21)
    \item \textbf{C10: Pitfall - Multiplying the Wrong Vector:} Multiplying $[x]_{\mathcal{C}}$ by $P_{\mathcal{C} \leftarrow \mathcal{B}}$ instead of its inverse. (Question: 22)
    \item \textbf{C11: Pitfall - Basis Order:} Assuming changing the order of a basis results in the identity matrix. (Question: 23)
    \item \textbf{C12: Pitfall - Inverse vs. Negative:} Negating the matrix instead of taking the inverse to reverse direction. (Question: 24)
    \item \textbf{C13: Pitfall - Dot Product Projection:} Using orthogonal projection formulas on non-orthogonal bases. (Question: 25)
    \item \textbf{C14: Advanced - Linear Transformations:} Changing the basis for a matrix $C = P^{-1}AP$. (Question: 26)
    \item \textbf{C15: Real-World Modeling (PCA):} Orthogonal matrices where $P^{-1} = P^T$. (Question: 27)
\end{itemize}

\newpage

\section{Problem Sets}

\subsection*{Type 1: Direct Matrix Construction \& Coordinate Transformation}
\textbf{Problem 1:} Let $\mathcal{B}=\{b_1, b_2\}$ and $\mathcal{C}=\{c_1, c_2\}$ be bases for a vector space $V$. Suppose $b_1 = 2c_1 - 3c_2$ and $b_2 = -c_1 + 5c_2$. Find the change-of-coordinates matrix from $\mathcal{B}$ to $\mathcal{C}$, denoted $P_{\mathcal{C} \leftarrow \mathcal{B}}$.

\vspace{0.5cm}
\textbf{Problem 2:} Using the bases from Problem 1, find the $\mathcal{C}$-coordinate vector $[x]_{\mathcal{C}}$ for the vector $x = 4b_1 + 2b_2$.

\vspace{0.5cm}
\textbf{Problem 30:} Let $\mathcal{B} = \{b_1, b_2, b_3\}$ and $\mathcal{C} = \{c_1, c_2, c_3\}$ be bases for $V$. Suppose $b_1 = c_1+c_2$, $b_2 = c_2+c_3$, and $b_3 = c_1+c_3$. Find $[x]_{\mathcal{C}}$ if $x = 2b_1 - b_2 + 3b_3$.

\subsection*{Type 2: Row Reduction Method in $\mathbb{R}^n$}
\textbf{Problem 3:} In $\mathbb{R}^2$, let $\mathcal{B} = \left\{ \begin{bmatrix} 1 \\ 2 \end{bmatrix}, \begin{bmatrix} 3 \\ 4 \end{bmatrix} \right\}$ and $\mathcal{C} = \left\{ \begin{bmatrix} 1 \\ -1 \end{bmatrix}, \begin{bmatrix} 0 \\ 1 \end{bmatrix} \right\}$. Find $P_{\mathcal{C} \leftarrow \mathcal{B}}$.

\vspace{0.5cm}
\textbf{Problem 4:} Using the matrices from Problem 3, find $[x]_{\mathcal{C}}$ if $[x]_{\mathcal{B}} = \begin{bmatrix} 2 \\ -1 \end{bmatrix}$.

\vspace{0.5cm}
\textbf{Problem 6:} In $\mathbb{R}^3$, let $\mathcal{B} = \left\{ \begin{bmatrix} 1 \\ 0 \\ 0 \end{bmatrix}, \begin{bmatrix} 1 \\ 1 \\ 0 \end{bmatrix}, \begin{bmatrix} 1 \\ 1 \\ 1 \end{bmatrix} \right\}$ and let $\mathcal{E}$ be the standard basis. Find $P_{\mathcal{E} \leftarrow \mathcal{B}}$.

\vspace{0.5cm}
\textbf{Problem 8:} In $\mathbb{R}^2$, let $\mathcal{B} = \left\{ \begin{bmatrix} 4 \\ 7 \end{bmatrix}, \begin{bmatrix} 1 \\ 2 \end{bmatrix} \right\}$ and $\mathcal{C} = \left\{ \begin{bmatrix} 1 \\ 1 \end{bmatrix}, \begin{bmatrix} 2 \\ 3 \end{bmatrix} \right\}$. Find $P_{\mathcal{C} \leftarrow \mathcal{B}}$.

\vspace{0.5cm}
\textbf{Problem 28:} In $\mathbb{R}^3$, let $\mathcal{B} = \left\{ \begin{bmatrix} 1 \\ 0 \\ 1 \end{bmatrix}, \begin{bmatrix} 0 \\ 1 \\ 1 \end{bmatrix}, \begin{bmatrix} 1 \\ 1 \\ 0 \end{bmatrix} \right\}$ and let $\mathcal{C}$ be the standard basis. Find $P_{\mathcal{C} \leftarrow \mathcal{B}}$.

\subsection*{Type 3: Invertibility \& Reversing Direction}
\textbf{Problem 5:} Using the result from Problem 3, find the change-of-coordinates matrix from $\mathcal{C}$ to $\mathcal{B}$, denoted $P_{\mathcal{B} \leftarrow \mathcal{C}}$.

\vspace{0.5cm}
\textbf{Problem 7:} Using the result from Problem 6, find $P_{\mathcal{B} \leftarrow \mathcal{E}}$.

\vspace{0.5cm}
\textbf{Problem 29:} Using the result from Problem 28, find $P_{\mathcal{B} \leftarrow \mathcal{C}}$.

\subsection*{Type 4: Abstract Vector Spaces (Polynomials \& Functions)}
\textbf{Problem 9:} In $\mathbb{P}_1$, find the change-of-coordinates matrix from the basis $\mathcal{B} = \{1+t, 1-t\}$ to the standard basis $\mathcal{C} = \{1, t\}$.

\vspace{0.5cm}
\textbf{Problem 10:} In $\mathbb{P}_2$, find the change-of-coordinates matrix from the basis $\mathcal{B} = \{1, 1+t, 1+t+t^2\}$ to the standard basis $\mathcal{C} = \{1, t, t^2\}$.

\vspace{0.5cm}
\textbf{Problem 11:} In $\mathbb{P}_2$, using the bases from Problem 10, find the $\mathcal{B}$-coordinate vector for the polynomial $p(t) = 2 - t + 3t^2$.

\vspace{0.5cm}
\textbf{Problem 12:} In $\mathbb{P}_2$, find the change-of-coordinates matrix from the basis $\mathcal{B} = \{1-t^2, t+t^2, 1+t\}$ to the standard basis $\mathcal{C} = \{1, t, t^2\}$.

\vspace{0.5cm}
\textbf{Problem 13:} Let $V = \text{span}\{\sin t, \cos t\}$. Find the change-of-coordinates matrix from $\mathcal{B} = \{\sin t + \cos t, 2\sin t - \cos t\}$ to $\mathcal{C} = \{\sin t, \cos t\}$.

\vspace{0.5cm}
\textbf{Problem 14:} Let $V = \text{span}\{e^t, te^t\}$. Find the change-of-coordinates matrix from $\mathcal{B} = \{e^t, te^t\}$ to $\mathcal{C} = \{e^t - te^t, e^t + te^t\}$.

\subsection*{Type 5: Conceptual True/False}
\textit{Mark each statement True or False and justify your answer.}

\vspace{0.5cm}
\textbf{Problem 15:} The columns of the change-of-coordinates matrix $P_{\mathcal{C} \leftarrow \mathcal{B}}$ are the $\mathcal{B}$-coordinate vectors of the vectors in $\mathcal{C}$.

\vspace{0.5cm}
\textbf{Problem 16:} If $V = \mathbb{R}^n$ and $\mathcal{C}$ is the standard basis, then $P_{\mathcal{C} \leftarrow \mathcal{B}}$ is simply the matrix whose columns are the vectors in $\mathcal{B}$.

\vspace{0.5cm}
\textbf{Problem 17:} The change-of-coordinates matrix $P_{\mathcal{C} \leftarrow \mathcal{B}}$ is always an invertible matrix.

\vspace{0.5cm}
\textbf{Problem 18:} If $\mathcal{B}$ and $\mathcal{C}$ are bases for $V$, and $x \in V$, then $[x]_{\mathcal{B}} = P_{\mathcal{C} \leftarrow \mathcal{B}} [x]_{\mathcal{C}}$.

\subsection*{Type 6: Theoretical Proofs \& Derivations}
\textbf{Problem 19:} Prove that for any two bases $\mathcal{B}$ and $\mathcal{C}$ of a vector space $V$, $P_{\mathcal{B} \leftarrow \mathcal{C}} P_{\mathcal{C} \leftarrow \mathcal{B}} = I$.

\vspace{0.5cm}
\textbf{Problem 20:} Let $T: V \to V$ be a linear transformation. Prove that if $A$ is the matrix of $T$ relative to basis $\mathcal{B}$ and $C$ is the matrix of $T$ relative to basis $\mathcal{C}$, then $C = P_{\mathcal{C} \leftarrow \mathcal{B}} A P_{\mathcal{B} \leftarrow \mathcal{C}}$.

\subsection*{Type 7: Advanced Diagnostic Testing (``Find the Flaw'')}
\textit{Identify the conceptual or computational flaw in the provided solutions and provide the correct solution.}

\vspace{0.5cm}
\textbf{Problem 21:} \textit{Goal:} Find $P_{\mathcal{C} \leftarrow \mathcal{B}}$ for bases in $\mathbb{R}^n$. \textit{Flawed Solution:} We set up the augmented matrix $[\mathcal{B} \mid \mathcal{C}]$ and row reduce it to $[I \mid P]$. The right side is $P_{\mathcal{C} \leftarrow \mathcal{B}}$. Find the flaw.

\vspace{0.5cm}
\textbf{Problem 22:} \textit{Goal:} Given $P_{\mathcal{C} \leftarrow \mathcal{B}} = \begin{bmatrix} 2 & 1 \\ 1 & 1 \end{bmatrix}$ and $[x]_{\mathcal{C}} = \begin{bmatrix} 3 \\ 2 \end{bmatrix}$, find $[x]_{\mathcal{B}}$. \textit{Flawed Solution:} We use the formula and multiply: $[x]_{\mathcal{B}} = \begin{bmatrix} 2 & 1 \\ 1 & 1 \end{bmatrix} \begin{bmatrix} 3 \\ 2 \end{bmatrix} = \begin{bmatrix} 8 \\ 5 \end{bmatrix}$. Find the flaw.

\vspace{0.5cm}
\textbf{Problem 23:} \textit{Goal:} Let $\mathcal{B} = \{v_1, v_2\}$ and $\mathcal{C} = \{v_2, v_1\}$. Find $P_{\mathcal{C} \leftarrow \mathcal{B}}$. \textit{Flawed Solution:} Because both bases contain the exact same vectors, the change of basis matrix is simply the identity matrix $I_2$. Find the flaw.

\vspace{0.5cm}
\textbf{Problem 24:} \textit{Goal:} If $P_{\mathcal{C} \leftarrow \mathcal{B}} = \begin{bmatrix} 3 & 5 \\ 1 & 2 \end{bmatrix}$, find $P_{\mathcal{B} \leftarrow \mathcal{C}}$. \textit{Flawed Solution:} To reverse the direction, we negate the matrix: $P_{\mathcal{B} \leftarrow \mathcal{C}} = \begin{bmatrix} -3 & -5 \\ -1 & -2 \end{bmatrix}$. Find the flaw.

\vspace{0.5cm}
\textbf{Problem 25:} \textit{Goal:} Let $\mathcal{B} = \left\{ \begin{bmatrix} 1 \\ 1 \end{bmatrix}, \begin{bmatrix} 1 \\ 2 \end{bmatrix} \right\}$. Find $[x]_{\mathcal{B}}$ for $x = \begin{bmatrix} 3 \\ 4 \end{bmatrix}$. \textit{Flawed Solution:} We use the dot product projection formula. $c_1 = \frac{x \cdot b_1}{b_1 \cdot b_1} = \frac{7}{2}$ and $c_2 = \frac{x \cdot b_2}{b_2 \cdot b_2} = \frac{11}{5}$. Find the flaw.

\subsection*{Type 8: Advanced Extensions (Linear Transformations \& PCA)}
\textbf{Problem 26:} Let $T(x) = \begin{bmatrix} 1 & 2 \\ 3 & 4 \end{bmatrix} x$ be a linear transformation in the standard basis $\mathcal{E}$. Find the matrix of $T$ relative to the basis $\mathcal{B} = \left\{ \begin{bmatrix} 1 \\ 1 \end{bmatrix}, \begin{bmatrix} 1 \\ -1 \end{bmatrix} \right\}$.

\vspace{0.5cm}
\textbf{Problem 27:} In Quantitative Finance (PCA), a covariance matrix $\Sigma$ is diagonalized as $\Sigma = P \Lambda P^{-1}$, where $P$ is the matrix of orthonormal eigenvectors (the eigenbasis $\mathcal{B}$). Explain why the change of basis matrix from the standard basis $\mathcal{E}$ to the eigenbasis $\mathcal{B}$ is given by $P^T$.

\newpage

\section{Detailed Solutions}

\textbf{Solution 1:} The change-of-coordinates matrix $P_{\mathcal{C} \leftarrow \mathcal{B}}$ is formed by placing the $\mathcal{C}$-coordinate vectors of the $\mathcal{B}$ basis vectors into columns. From the given equations:
$[b_1]_{\mathcal{C}} = \begin{bmatrix} 2 \\ -3 \end{bmatrix}$ and $[b_2]_{\mathcal{C}} = \begin{bmatrix} -1 \\ 5 \end{bmatrix}$.
Therefore, $P_{\mathcal{C} \leftarrow \mathcal{B}} = \begin{bmatrix} 2 & -1 \\ -3 & 5 \end{bmatrix}$.

\vspace{0.5cm}
\textbf{Solution 2:} We are given $x = 4b_1 + 2b_2$, which means $[x]_{\mathcal{B}} = \begin{bmatrix} 4 \\ 2 \end{bmatrix}$.
Using the change of coordinates equation $[x]_{\mathcal{C}} = P_{\mathcal{C} \leftarrow \mathcal{B}}[x]_{\mathcal{B}}$:
$[x]_{\mathcal{C}} = \begin{bmatrix} 2 & -1 \\ -3 & 5 \end{bmatrix} \begin{bmatrix} 4 \\ 2 \end{bmatrix} = \begin{bmatrix} (2)(4) + (-1)(2) \\ (-3)(4) + (5)(2) \end{bmatrix} = \begin{bmatrix} 8 - 2 \\ -12 + 10 \end{bmatrix} = \begin{bmatrix} 6 \\ -2 \end{bmatrix}$.

\vspace{0.5cm}
\textbf{Solution 30:} First, construct $P_{\mathcal{C} \leftarrow \mathcal{B}}$ by extracting the coefficients:
$[b_1]_{\mathcal{C}} = \begin{bmatrix} 1 \\ 1 \\ 0 \end{bmatrix}$, $[b_2]_{\mathcal{C}} = \begin{bmatrix} 0 \\ 1 \\ 1 \end{bmatrix}$, $[b_3]_{\mathcal{C}} = \begin{bmatrix} 1 \\ 0 \\ 1 \end{bmatrix}$.
$P_{\mathcal{C} \leftarrow \mathcal{B}} = \begin{bmatrix} 1 & 0 & 1 \\ 1 & 1 & 0 \\ 0 & 1 & 1 \end{bmatrix}$.
We are given $x = 2b_1 - b_2 + 3b_3$, so $[x]_{\mathcal{B}} = \begin{bmatrix} 2 \\ -1 \\ 3 \end{bmatrix}$.
$[x]_{\mathcal{C}} = P_{\mathcal{C} \leftarrow \mathcal{B}}[x]_{\mathcal{B}} = \begin{bmatrix} 1 & 0 & 1 \\ 1 & 1 & 0 \\ 0 & 1 & 1 \end{bmatrix} \begin{bmatrix} 2 \\ -1 \\ 3 \end{bmatrix} = \begin{bmatrix} 2 + 0 + 3 \\ 2 - 1 + 0 \\ 0 - 1 + 3 \end{bmatrix} = \begin{bmatrix} 5 \\ 1 \\ 2 \end{bmatrix}$.

\vspace{0.5cm}
\textbf{Solution 3:} We use the augmented matrix method $[\mathcal{C} \mid \mathcal{B}] \sim [I \mid P_{\mathcal{C} \leftarrow \mathcal{B}}]$.
$[\mathcal{C} \mid \mathcal{B}] = \begin{bmatrix} 1 & 0 & 1 & 3 \\ -1 & 1 & 2 & 4 \end{bmatrix}$.
Row reduce: $R_2 \leftarrow R_2 + R_1$:
$\begin{bmatrix} 1 & 0 & 1 & 3 \\ 0 & 1 & 3 & 7 \end{bmatrix}$.
The right half is $P_{\mathcal{C} \leftarrow \mathcal{B}} = \begin{bmatrix} 1 & 3 \\ 3 & 7 \end{bmatrix}$.

\vspace{0.5cm}
\textbf{Solution 4:} Using $P_{\mathcal{C} \leftarrow \mathcal{B}}$ from Problem 3:
$[x]_{\mathcal{C}} = P_{\mathcal{C} \leftarrow \mathcal{B}}[x]_{\mathcal{B}} = \begin{bmatrix} 1 & 3 \\ 3 & 7 \end{bmatrix} \begin{bmatrix} 2 \\ -1 \end{bmatrix} = \begin{bmatrix} 2 - 3 \\ 6 - 7 \end{bmatrix} = \begin{bmatrix} -1 \\ -1 \end{bmatrix}$.

\vspace{0.5cm}
\textbf{Solution 6:} Since $\mathcal{E}$ is the standard basis, we can use the \textbf{Standard Basis Shortcut}. $P_{\mathcal{E} \leftarrow \mathcal{B}}$ is simply the matrix whose columns are the vectors of $\mathcal{B}$.
$P_{\mathcal{E} \leftarrow \mathcal{B}} = \begin{bmatrix} 1 & 1 & 1 \\ 0 & 1 & 1 \\ 0 & 0 & 1 \end{bmatrix}$.

\vspace{0.5cm}
\textbf{Solution 8:} Set up $[\mathcal{C} \mid \mathcal{B}]$:
$\begin{bmatrix} 1 & 2 & 4 & 1 \\ 1 & 3 & 7 & 2 \end{bmatrix}$.
$R_2 \leftarrow R_2 - R_1$: $\begin{bmatrix} 1 & 2 & 4 & 1 \\ 0 & 1 & 3 & 1 \end{bmatrix}$.
$R_1 \leftarrow R_1 - 2R_2$: $\begin{bmatrix} 1 & 0 & -2 & -1 \\ 0 & 1 & 3 & 1 \end{bmatrix}$.
$P_{\mathcal{C} \leftarrow \mathcal{B}} = \begin{bmatrix} -2 & -1 \\ 3 & 1 \end{bmatrix}$.

\vspace{0.5cm}
\textbf{Solution 28:} Since $\mathcal{C}$ is the standard basis, $P_{\mathcal{C} \leftarrow \mathcal{B}}$ is just the vectors of $\mathcal{B}$ as columns.
$P_{\mathcal{C} \leftarrow \mathcal{B}} = \begin{bmatrix} 1 & 0 & 1 \\ 0 & 1 & 1 \\ 1 & 1 & 0 \end{bmatrix}$.

\vspace{0.5cm}
\textbf{Solution 5:} A fundamental theorem states $P_{\mathcal{B} \leftarrow \mathcal{C}} = (P_{\mathcal{C} \leftarrow \mathcal{B}})^{-1}$.
From Problem 3, $P = \begin{bmatrix} 1 & 3 \\ 3 & 7 \end{bmatrix}$.
Determinant $= (1)(7) - (3)(3) = -2$.
$P^{-1} = \frac{1}{-2} \begin{bmatrix} 7 & -3 \\ -3 & 1 \end{bmatrix} = \begin{bmatrix} -3.5 & 1.5 \\ 1.5 & -0.5 \end{bmatrix}$.

\vspace{0.5cm}
\textbf{Solution 7:} $P_{\mathcal{B} \leftarrow \mathcal{E}} = (P_{\mathcal{E} \leftarrow \mathcal{B}})^{-1}$.
From Problem 6, $P = \begin{bmatrix} 1 & 1 & 1 \\ 0 & 1 & 1 \\ 0 & 0 & 1 \end{bmatrix}$.
Row reducing $[P \mid I]$ yields the inverse:
$P^{-1} = \begin{bmatrix} 1 & -1 & 0 \\ 0 & 1 & -1 \\ 0 & 0 & 1 \end{bmatrix}$.

\vspace{0.5cm}
\textbf{Solution 29:} $P_{\mathcal{B} \leftarrow \mathcal{C}} = (P_{\mathcal{C} \leftarrow \mathcal{B}})^{-1}$.
From Problem 28, $P = \begin{bmatrix} 1 & 0 & 1 \\ 0 & 1 & 1 \\ 1 & 1 & 0 \end{bmatrix}$.
Determinant $= 1(0 - 1) - 0 + 1(0 - 1) = -2$.
Using the adjugate method or row reduction, the inverse is:
$P^{-1} = \begin{bmatrix} 0.5 & -0.5 & 0.5 \\ -0.5 & 0.5 & 0.5 \\ 0.5 & 0.5 & -0.5 \end{bmatrix}$.

\vspace{0.5cm}
\textbf{Solution 9:} Map the abstract vectors to $\mathbb{R}^2$ using the standard basis $\mathcal{C} = \{1, t\}$.
$[1+t]_{\mathcal{C}} = \begin{bmatrix} 1 \\ 1 \end{bmatrix}$, $[1-t]_{\mathcal{C}} = \begin{bmatrix} 1 \\ -1 \end{bmatrix}$.
$P_{\mathcal{C} \leftarrow \mathcal{B}} = \begin{bmatrix} 1 & 1 \\ 1 & -1 \end{bmatrix}$.

\vspace{0.5cm}
\textbf{Solution 10:} Map to $\mathbb{R}^3$ using $\mathcal{C} = \{1, t, t^2\}$.
$[1]_{\mathcal{C}} = \begin{bmatrix} 1 \\ 0 \\ 0 \end{bmatrix}$, $[1+t]_{\mathcal{C}} = \begin{bmatrix} 1 \\ 1 \\ 0 \end{bmatrix}$, $[1+t+t^2]_{\mathcal{C}} = \begin{bmatrix} 1 \\ 1 \\ 1 \end{bmatrix}$.
$P_{\mathcal{C} \leftarrow \mathcal{B}} = \begin{bmatrix} 1 & 1 & 1 \\ 0 & 1 & 1 \\ 0 & 0 & 1 \end{bmatrix}$.

\vspace{0.5cm}
\textbf{Solution 11:} We need $[p]_{\mathcal{B}}$. We know $[p]_{\mathcal{C}} = \begin{bmatrix} 2 \\ -1 \\ 3 \end{bmatrix}$.
We must solve $P_{\mathcal{C} \leftarrow \mathcal{B}}[p]_{\mathcal{B}} = [p]_{\mathcal{C}}$, so $[p]_{\mathcal{B}} = (P_{\mathcal{C} \leftarrow \mathcal{B}})^{-1}[p]_{\mathcal{C}}$.
From Solution 7, $P^{-1} = \begin{bmatrix} 1 & -1 & 0 \\ 0 & 1 & -1 \\ 0 & 0 & 1 \end{bmatrix}$.
$[p]_{\mathcal{B}} = \begin{bmatrix} 1 & -1 & 0 \\ 0 & 1 & -1 \\ 0 & 0 & 1 \end{bmatrix} \begin{bmatrix} 2 \\ -1 \\ 3 \end{bmatrix} = \begin{bmatrix} 2 - (-1) \\ -1 - 3 \\ 3 \end{bmatrix} = \begin{bmatrix} 3 \\ -4 \\ 3 \end{bmatrix}$.

\vspace{0.5cm}
\textbf{Solution 12:} Map to $\mathbb{R}^3$ using $\mathcal{C} = \{1, t, t^2\}$.
$[1-t^2]_{\mathcal{C}} = \begin{bmatrix} 1 \\ 0 \\ -1 \end{bmatrix}$, $[t+t^2]_{\mathcal{C}} = \begin{bmatrix} 0 \\ 1 \\ 1 \end{bmatrix}$, $[1+t]_{\mathcal{C}} = \begin{bmatrix} 1 \\ 1 \\ 0 \end{bmatrix}$.
$P_{\mathcal{C} \leftarrow \mathcal{B}} = \begin{bmatrix} 1 & 0 & 1 \\ 0 & 1 & 1 \\ -1 & 1 & 0 \end{bmatrix}$.

\vspace{0.5cm}
\textbf{Solution 13:} Map the functions to $\mathbb{R}^2$ using $\mathcal{C} = \{\sin t, \cos t\}$.
$[\sin t + \cos t]_{\mathcal{C}} = \begin{bmatrix} 1 \\ 1 \end{bmatrix}$, $[2\sin t - \cos t]_{\mathcal{C}} = \begin{bmatrix} 2 \\ -1 \end{bmatrix}$.
$P_{\mathcal{C} \leftarrow \mathcal{B}} = \begin{bmatrix} 1 & 2 \\ 1 & -1 \end{bmatrix}$.

\vspace{0.5cm}
\textbf{Solution 14:} Map to $\mathbb{R}^2$ using the standard basis for this space, let's call it $\mathcal{E} = \{e^t, te^t\}$.
Then $\mathcal{B}$ is the standard basis, and $\mathcal{C}$ is nonstandard.
We want $P_{\mathcal{C} \leftarrow \mathcal{B}}$. Set up $[\mathcal{C} \mid \mathcal{B}]$:
$\begin{bmatrix} 1 & 1 & 1 & 0 \\ -1 & 1 & 0 & 1 \end{bmatrix}$.
$R_2 \leftarrow R_2 + R_1$: $\begin{bmatrix} 1 & 1 & 1 & 0 \\ 0 & 2 & 1 & 1 \end{bmatrix}$.
$R_2 \leftarrow R_2 / 2$: $\begin{bmatrix} 1 & 1 & 1 & 0 \\ 0 & 1 & 0.5 & 0.5 \end{bmatrix}$.
$R_1 \leftarrow R_1 - R_2$: $\begin{bmatrix} 1 & 0 & 0.5 & -0.5 \\ 0 & 1 & 0.5 & 0.5 \end{bmatrix}$.
$P_{\mathcal{C} \leftarrow \mathcal{B}} = \begin{bmatrix} 0.5 & -0.5 \\ 0.5 & 0.5 \end{bmatrix}$.

\vspace{0.5cm}
\textbf{Solution 15: False.} The columns of $P_{\mathcal{C} \leftarrow \mathcal{B}}$ are the $\mathcal{C}$-coordinate vectors of the vectors in $\mathcal{B}$, not the other way around. The notation means we are converting from $\mathcal{B}$ to $\mathcal{C}$, so we evaluate $\mathcal{B}$ in the language of $\mathcal{C}$.

\vspace{0.5cm}
\textbf{Solution 16: True.} If $\mathcal{C}$ is the standard basis $\mathcal{E}$, then the $\mathcal{E}$-coordinate vector of any vector $x$ in $\mathbb{R}^n$ is just $x$ itself. Thus, the columns are exactly the vectors of $\mathcal{B}$.

\vspace{0.5cm}
\textbf{Solution 17: True.} The columns of $P_{\mathcal{C} \leftarrow \mathcal{B}}$ are the coordinate vectors of a basis. Because basis vectors are linearly independent, their coordinate vectors are linearly independent. A square matrix with linearly independent columns is invertible.

\vspace{0.5cm}
\textbf{Solution 18: False.} The correct equation is $[x]_{\mathcal{C}} = P_{\mathcal{C} \leftarrow \mathcal{B}} [x]_{\mathcal{B}}$. The matrix $P_{\mathcal{C} \leftarrow \mathcal{B}}$ consumes a $\mathcal{B}$-vector to produce a $\mathcal{C}$-vector.

\vspace{0.5cm}
\textbf{Solution 19:} By definition, for any vector $x \in V$, $[x]_{\mathcal{C}} = P_{\mathcal{C} \leftarrow \mathcal{B}} [x]_{\mathcal{B}}$.
Applying the reverse transformation, $[x]_{\mathcal{B}} = P_{\mathcal{B} \leftarrow \mathcal{C}} [x]_{\mathcal{C}}$.
Substituting the first equation into the second gives:
$[x]_{\mathcal{B}} = P_{\mathcal{B} \leftarrow \mathcal{C}} (P_{\mathcal{C} \leftarrow \mathcal{B}} [x]_{\mathcal{B}}) = (P_{\mathcal{B} \leftarrow \mathcal{C}} P_{\mathcal{C} \leftarrow \mathcal{B}}) [x]_{\mathcal{B}}$.
Since this holds for all coordinate vectors $[x]_{\mathcal{B}}$ in $\mathbb{R}^n$, the matrix product $P_{\mathcal{B} \leftarrow \mathcal{C}} P_{\mathcal{C} \leftarrow \mathcal{B}}$ must be the identity matrix $I$.

\vspace{0.5cm}
\textbf{Solution 20:} By definition of a matrix for a linear transformation, $[T(x)]_{\mathcal{C}} = C[x]_{\mathcal{C}}$.
We can also find $[T(x)]_{\mathcal{C}}$ by converting to basis $\mathcal{B}$, applying $A$, and converting back:
$[T(x)]_{\mathcal{C}} = P_{\mathcal{C} \leftarrow \mathcal{B}} [T(x)]_{\mathcal{B}} = P_{\mathcal{C} \leftarrow \mathcal{B}} (A [x]_{\mathcal{B}}) = P_{\mathcal{C} \leftarrow \mathcal{B}} A (P_{\mathcal{B} \leftarrow \mathcal{C}} [x]_{\mathcal{C}})$.
Thus, $C[x]_{\mathcal{C}} = (P_{\mathcal{C} \leftarrow \mathcal{B}} A P_{\mathcal{B} \leftarrow \mathcal{C}}) [x]_{\mathcal{C}}$ for all $x$. Therefore, $C = P_{\mathcal{C} \leftarrow \mathcal{B}} A P_{\mathcal{B} \leftarrow \mathcal{C}}$.

\vspace{0.5cm}
\textbf{Solution 21: Flaw:} The student set up the augmented matrix backwards. To find $P_{\mathcal{C} \leftarrow \mathcal{B}}$, the ``target'' basis $\mathcal{C}$ must go on the left, and the ``source'' basis $\mathcal{B}$ on the right. \textbf{Correct Solution:} Set up $[\mathcal{C} \mid \mathcal{B}]$ and row reduce to $[I \mid P_{\mathcal{C} \leftarrow \mathcal{B}}]$.

\vspace{0.5cm}
\textbf{Solution 22: Flaw:} The student multiplied the $\mathcal{C}$-coordinate vector by $P_{\mathcal{C} \leftarrow \mathcal{B}}$, which requires a $\mathcal{B}$-coordinate vector as input. \textbf{Correct Solution:} We must use the inverse matrix: $[x]_{\mathcal{B}} = (P_{\mathcal{C} \leftarrow \mathcal{B}})^{-1} [x]_{\mathcal{C}}$.
$(P_{\mathcal{C} \leftarrow \mathcal{B}})^{-1} = \frac{1}{2-1} \begin{bmatrix} 1 & -1 \\ -1 & 2 \end{bmatrix} = \begin{bmatrix} 1 & -1 \\ -1 & 2 \end{bmatrix}$.
$[x]_{\mathcal{B}} = \begin{bmatrix} 1 & -1 \\ -1 & 2 \end{bmatrix} \begin{bmatrix} 3 \\ 2 \end{bmatrix} = \begin{bmatrix} 1 \\ 1 \end{bmatrix}$.

\vspace{0.5cm}
\textbf{Solution 23: Flaw:} The order of vectors in a basis strictly dictates the coordinate mapping. Changing the order changes the coordinate vectors. \textbf{Correct Solution:}
$[b_1]_{\mathcal{C}} = [v_1]_{\mathcal{C}} = \begin{bmatrix} 0 \\ 1 \end{bmatrix}$ (since $v_1 = 0v_2 + 1v_1$).
$[b_2]_{\mathcal{C}} = [v_2]_{\mathcal{C}} = \begin{bmatrix} 1 \\ 0 \end{bmatrix}$.
$P_{\mathcal{C} \leftarrow \mathcal{B}} = \begin{bmatrix} 0 & 1 \\ 1 & 0 \end{bmatrix}$.

\vspace{0.5cm}
\textbf{Solution 24: Flaw:} Change of basis is a linear transformation (multiplication), not a translation (addition). The reverse operation is the matrix inverse, not the negative. \textbf{Correct Solution:}
$P_{\mathcal{B} \leftarrow \mathcal{C}} = (P_{\mathcal{C} \leftarrow \mathcal{B}})^{-1} = \frac{1}{6-5} \begin{bmatrix} 2 & -5 \\ -1 & 3 \end{bmatrix} = \begin{bmatrix} 2 & -5 \\ -1 & 3 \end{bmatrix}$.

\vspace{0.5cm}
\textbf{Solution 25: Flaw:} The dot product projection formula $c_i = \frac{x \cdot b_i}{b_i \cdot b_i}$ only works if the basis $\mathcal{B}$ is an orthogonal basis. This basis is not orthogonal ($b_1 \cdot b_2 = 1(1) + 1(2) = 3 \neq 0$). \textbf{Correct Solution:} Solve the system $c_1b_1 + c_2b_2 = x$.
$\begin{bmatrix} 1 & 1 \\ 1 & 2 \end{bmatrix} \begin{bmatrix} c_1 \\ c_2 \end{bmatrix} = \begin{bmatrix} 3 \\ 4 \end{bmatrix}$.
$R_2 \leftarrow R_2 - R_1 \implies c_2 = 1$.
$c_1 + 1 = 3 \implies c_1 = 2$.
$[x]_{\mathcal{B}} = \begin{bmatrix} 2 \\ 1 \end{bmatrix}$.

\vspace{0.5cm}
\textbf{Solution 26:} We use the formula $B = P_{\mathcal{B} \leftarrow \mathcal{E}} A P_{\mathcal{E} \leftarrow \mathcal{B}} = P^{-1} A P$.
$P = P_{\mathcal{E} \leftarrow \mathcal{B}} = \begin{bmatrix} 1 & 1 \\ 1 & -1 \end{bmatrix}$.
$P^{-1} = \frac{1}{-2} \begin{bmatrix} -1 & -1 \\ -1 & 1 \end{bmatrix} = \begin{bmatrix} 0.5 & 0.5 \\ 0.5 & -0.5 \end{bmatrix}$.
$AP = \begin{bmatrix} 1 & 2 \\ 3 & 4 \end{bmatrix} \begin{bmatrix} 1 & 1 \\ 1 & -1 \end{bmatrix} = \begin{bmatrix} 3 & -1 \\ 7 & -1 \end{bmatrix}$.
$B = P^{-1} (AP) = \begin{bmatrix} 0.5 & 0.5 \\ 0.5 & -0.5 \end{bmatrix} \begin{bmatrix} 3 & -1 \\ 7 & -1 \end{bmatrix} = \begin{bmatrix} 1.5 + 3.5 & -0.5 - 0.5 \\ 1.5 - 3.5 & -0.5 + 0.5 \end{bmatrix} = \begin{bmatrix} 5 & -1 \\ -2 & 0 \end{bmatrix}$.

\vspace{0.5cm}
\textbf{Solution 27:} The matrix $P$ contains the orthonormal eigenvectors as its columns. By definition, this makes $P$ the change-of-coordinates matrix from the eigenbasis $\mathcal{B}$ to the standard basis $\mathcal{E}$, so $P = P_{\mathcal{E} \leftarrow \mathcal{B}}$.
Because the columns of $P$ are orthonormal, $P$ is an orthogonal matrix, which has the special property that its inverse is its transpose ($P^{-1} = P^T$).
The change of basis from standard to eigenbasis is $P_{\mathcal{B} \leftarrow \mathcal{E}} = (P_{\mathcal{E} \leftarrow \mathcal{B}})^{-1} = P^{-1}$.
Therefore, $P_{\mathcal{B} \leftarrow \mathcal{E}} = P^T$.

\newpage

\section{Concept Checklist \& Question Mapping (Appendix)}
\begin{itemize}
    \item \textbf{C1: Direct Matrix Construction:} Extracting coefficients directly from linear combinations to form $P_{\mathcal{C} \leftarrow \mathcal{B}}$. (Questions: 1, 2, 30)
    \item \textbf{C2: Row Reduction Method in $\mathbb{R}^n$:} Using the augmented matrix $[\mathcal{C} \mid \mathcal{B}] \sim [I \mid P_{\mathcal{C} \leftarrow \mathcal{B}}]$. (Questions: 3, 4, 6, 8, 28)
    \item \textbf{C3: Coordinate Transformation:} Using the equation $[x]_{\mathcal{C}} = P_{\mathcal{C} \leftarrow \mathcal{B}}[x]_{\mathcal{B}}$. (Questions: 2, 4, 11, 30)
    \item \textbf{C4: Invertibility \& Reversing Direction:} Using $P_{\mathcal{B} \leftarrow \mathcal{C}} = (P_{\mathcal{C} \leftarrow \mathcal{B}})^{-1}$. (Questions: 5, 7, 29)
    \item \textbf{C5: Abstract Vector Spaces (Polynomials):} Isomorphic mapping of $\mathbb{P}_n$ to $\mathbb{R}^{n+1}$. (Questions: 9, 10, 11, 12)
    \item \textbf{C6: Abstract Vector Spaces (Functions/Trig):} Mapping continuous functions to discrete matrices. (Questions: 13, 14)
    \item \textbf{C7: Conceptual True/False:} Testing definitions, linear independence, and matrix properties. (Questions: 15, 16, 17, 18)
    \item \textbf{C8: Theoretical Proofs \& Derivations:} Relying on definitions of linear transformations. (Questions: 19, 20)
    \item \textbf{C9: Pitfall - Reversing the Direction:} Setting up the augmented matrix backwards as $[\mathcal{B} \mid \mathcal{C}]$. (Question: 21)
    \item \textbf{C10: Pitfall - Multiplying the Wrong Vector:} Multiplying $[x]_{\mathcal{C}}$ by $P_{\mathcal{C} \leftarrow \mathcal{B}}$ instead of its inverse. (Question: 22)
    \item \textbf{C11: Pitfall - Basis Order:} Assuming changing the order of a basis results in the identity matrix. (Question: 23)
    \item \textbf{C12: Pitfall - Inverse vs. Negative:} Negating the matrix instead of taking the inverse to reverse direction. (Question: 24)
    \item \textbf{C13: Pitfall - Dot Product Projection:} Using orthogonal projection formulas on non-orthogonal bases. (Question: 25)
    \item \textbf{C14: Advanced - Linear Transformations:} Changing the basis for a matrix $C = P^{-1}AP$. (Question: 26)
    \item \textbf{C15: Real-World Modeling (PCA):} Orthogonal matrices where $P^{-1} = P^T$. (Question: 27)
\end{itemize}

\end{document}